% 文件开头添加条件判断
\newif\ifstandalone
\standalonetrue  % 默认独立编译模式

% 检测是否在主文档中
\ifdefined\maindoc
  \standalonefalse
\fi

\ifstandalone
  % C-附录C-常见机器人动力学.tex
  \documentclass[UTF8]{ctexart}

  \usepackage{amsmath,amssymb,amsfonts}
  \usepackage{geometry}
  \geometry{a4paper,margin=2.5cm}
  \usepackage{graphicx}
  \usepackage{float}
\fi

\ifstandalone
  \begin{document}
  \section*{附录C}
  \addcontentsline{toc}{section}{附录C 常见机器人机械臂的动力学}
\fi

% 当在主文档中使用时
\ifstandalone
\else
  \section*{附录C}
  \addcontentsline{toc}{section}{附录C 常见机器人机械臂的动力学}
\fi

\subsection*{常见机器人机械臂的动力学}

在本附录中，我们给出一些常见机器人机械臂的动力学。假设机器人动力学由下式给出：

\begin{equation}
\tau = M(q)\ddot{q} + V_m(q, \dot{q})\dot{q} + G(q) = M(q)\ddot{q} + N(q, \dot{q})
\tag{C.1.1}
\end{equation}

其中矩阵 $M(q)$ 是对称且正定的，其元素为 $m_{ij}(q)$，即

\begin{align*}
M(q) &= [m_{ij}(q)]_{n \times n}, \quad i,j = 1, \ldots, n \\
m_{ij}(q) &= m_{ji}(q)
\end{align*}

而 $N(q, \dot{q})$ 是一个 $n \times 1$ 向量，其元素为 $n_i$，即

\begin{equation*}
N(q, \dot{q}) = [n_i(q, \dot{q})]_{n \times 1}
\end{equation*}

特别要注意，在 $n_i$ 的表达式中通过重力常数 $g = 9.8$ 米/秒$^2$ 来识别重力项。我们还将采用以下符号表示：

\begin{itemize}
\item 连杆 $i$ 的长度为 $L_i$，单位为米
\item 连杆 $i$ 的质量为 $m_i$，单位为千克
\item 连杆 $i$ 关于轴 $u$ 的质量惯性矩为 $I_{uui}$，单位为 kg-m-m
\item $S_i = \sin q_i$，$C_i = \cos q_i$
\item $S_{ij} = \sin(q_i + q_j)$，$C_{ij} = \cos(q_i + q_j)$
\item $S_{ijk} = \sin(q_i + q_j + q_k)$，$C_{ijk} = \cos(q_i + q_j + q_k)$
\item $SS_i = \sin^2 q_i$，$CC_i = \cos^2 q_i$，$CS_i = \cos q_i \sin q_i$
\item $SS_{ij} = \sin^2(q_i + q_j)$
\end{itemize}

\subsection*{C.1 SCARA 机械臂}

我们考虑的第一个机器人是通用 SCARA 构型机器人，如图 C.1.1 所示。这些方程适用于 AdeptOne 和 AdeptTwo 机器人。动力学包括前四个自由度，以符号形式给出如下：

\begin{align*}
m_{11} &= (I_{zz1} + I_{zz2} + I_{zz3} + I_{zz4}) + (\frac{m_1}{4} + 2m_3 + m_4)L_1^2 \notag\\
&\quad + (\frac{m_2}{2} + 3m_3 + m_4)L_1L_2C_2 + (\frac{m_2}{4} + m_3 + m_4)L_2^2 \\
m_{12} &= (I_{zz2} + I_{zz3} + I_{zz4}) + (\frac{m_2}{4} + m_3 + m_4)L_2^2 \notag\\
&\quad + (2m_3 + m_4)L_1L_2C_2 \\
m_{13} &= 0 \\
m_{14} &= I_{zz4} \\
m_{22} &= (I_{zz2} + I_{zz3} + I_{zz4}) + (\frac{m_2}{4} + m_3 + m_4)L_2^2 \\
m_{23} &= 0 \\
m_{24} &= I_{zz4} \\
m_{33} &= m_3 + m_4 \\
m_{34} &= 0 \\
m_{44} &= I_{zz4}
\end{align*}

\begin{align*}
n_1 &= (\frac{m_2}{2} - m_3)L_1L_2\dot{q}_1^2 - (4m_3 + 2m_4)L_1L_2\dot{q}_1\dot{q}_2 \\
n_2 &= (\frac{m_2}{2} + m_3 + m_4)L_1L_2S_2\dot{q}_1^2 \\
n_3 &= (m_3 + m_4)g \\
n_4 &= 0
\end{align*}

\begin{figure}[H]
\centering
\fbox{\parbox{0.8\textwidth}{\centering 图 C.1.1：SCARA 机械臂}}
\caption{SCARA 机械臂示意图}
\end{figure}

\subsection*{C.2 Stanford 机械臂}

如图 C.2.1 所示的 Stanford 机械臂具有以下动力学方程 [Bejczy 1974], [Paul 1981]：

\begin{align*}
M_{11} &= 1.316 - 1.056108d_2 + 11.48d_2^2 + 2.51S_2^2 \notag\\
&\quad - 5.47995S_2^2d_3 + 6.47S_2^2d_3^2 + 0.23S_2^2d_3C_5 \\
m_{12} &= -6.47C_2d_2d_3 \\
m_{13} &= -6.47S_2d_2 \\
m_{14} &= 0 \\
m_{15} &= 0 \\
m_{16} &= 0 \\
m_{22} &= 4.721 - 5.47995d_3 + 6.47d_3^2 + 0.23d_3C_5 \\
m_{23} &= 0 \\
m_{24} &= 0 \\
m_{25} &= 0 \\
m_{26} &= 0 \\
m_{33} &= 7.252 \\
m_{34} &= 0 \\
m_{35} &= 0 \\
m_{36} &= 0 \\
m_{44} &= 0.107 + 0.203S_5^2 \\
m_{45} &= 0 \\
m_{46} &= 0 \\
m_{55} &= 0.113 \\
m_{56} &= 0 \\
m_{66} &= 0.0203
\end{align*}

\begin{align*}
n_1 &= 0 \\
n_2 &= -[2.734S_2 + 6.47S_2d_3 + 0.115(S_2C_5 + C_2C_4S_5)]g \\
n_3 &= 6.47C_2g \\
n_4 &= 0.115(S_2S_4S_5)g \\
n_5 &= 0.115(S_2C_4C_5 - C_2S_5)g \\
n_6 &= 0
\end{align*}

\begin{figure}[H]
\centering
\fbox{\parbox{0.8\textwidth}{\centering 图 C.2.1：Stanford 机械臂}}
\caption{Stanford 机械臂示意图}
\end{figure}

\subsection*{C.3 PUMA 560 机械臂}

PUMA 560 如图 C.3.1 所示。对于这种特定的结构，可以进行许多简化以获得以下动力学方程，这些方程出现在 [Armstrong et al. 1986] 中。

\begin{align*}
m_{11} &= 2.57 + 1.38CC_2 + 0.3SS_{23} + 0.744C_2S_{23} \\
m_{12} &= 0.69S_2 - 0.134C_{23} - 0.0238C_2 \\
m_{13} &= -0.134C_{23} - 0.00397S_{23} \\
m_{14} &= 0 \\
m_{15} &= 0 \\
m_{16} &= 0 \\
m_{22} &= 6.79 + 0.744S_3 \\
m_{23} &= 0.333 + 0.372S_3 - 0.011C_3 \\
m_{24} &= 0 \\
m_{25} &= 0 \\
m_{26} &= 0 \\
m_{33} &= 1.16 \\
m_{34} &= -0.00125S_4S_5 \\
m_{35} &= -0.00125C_4C_5 \\
m_{36} &= 0 \\
m_{44} &= 0.2 \\
m_{45} &= 0 \\
m_{46} &= 0 \\
m_{55} &= 0.18 \\
m_{56} &= 0 \\
m_{66} &= 0.19
\end{align*}

\begin{align*}
n_1 &= [0.69C_2 + 0.134S_{23} - 0.0238S_2]\dot{q}_2^2 + [0.1335S_{23} - 0.00379C_{23}]\dot{q}_3^2 \\
&\quad + [-2.76SC_2 + 0.744C_{223} + 0.6SC_{23} - 0.0213(1 - 2SS_{23})]\dot{q}_1\dot{q}_2 \\
&\quad + [0.744C_2C_{23} + 0.6SC_{23} + 0.022C_2S_{23} - 0.0213(1 - 2SS_{23})]\dot{q}_1\dot{q}_3 \\
&\quad + [-0.0025SC_{23}S_4S_5 - 0.00086C_4S_5 + 0.00248C_2C_{23}S_4S_5]\dot{q}_1\dot{q}_4 \\
&\quad + [-0.0025(SS_{23}S_5 - SC_{23}C_4C_5) - 0.00248C_2(S_{23}S_5 - C_{23}C_4C_5) \\
&\quad + 0.000864S_4C_5]\dot{q}_1\dot{q}_5 + [0.267S_{23} - 0.00758C_{23}]\dot{q}_2\dot{q}_3
\end{align*}

\begin{align*}
n_2 &= -\frac{1}{2}[-2.76SC_2 + 0.744C_{223} + 0.6SC_{23} - 0.0213(1 - 2SS_{23})]\dot{q}_1^2 \\
&\quad + \frac{1}{2}[0.022S_3 + 0.744C_3]\dot{q}_3^2 \\
&\quad + [0.00164S_{23} - 0.0025C_{23}C_4S_5 + 0.00248S_2C_4S_5 + 0.00003S_{23}(1 - 2SS_4)]\dot{q}_1\dot{q}_4 \\
&\quad + [-0.00215C_{23}S_4C_5 + 0.00248S_2S_4C_5 - 0.000642C_{23}S_4]\dot{q}_1\dot{q}_5 \\
&\quad + [0.022S_3 + 0.744C_3]\dot{q}_2\dot{q}_3 - [0.00248C_3S_4S_5]\dot{q}_2\dot{q}_4 \\
&\quad + [-0.0025S_5 + 0.00248(C_3C_4C_5 - S_3S_5)]\dot{q}_2\dot{q}_5 \\
&\quad - [0.00248C_3S_4S_5]\dot{q}_3\dot{q}_4 \\
&\quad + [-0.0025S_5 + 0.00248(C_3C_4C_5 - S_3S_5)]\dot{q}_3\dot{q}_5 \\
&\quad - 37.2C_2 - 8.4S_{23} + 1.02S_2
\end{align*}

\begin{align*}
n_3 &= -\frac{1}{2}[-2.76SC_2 + 0.744C_{223} + 0.6SC_{23} - 0.0213(1 - 2SS_{23})]\dot{q}_1^2 \\
&\quad - \frac{1}{2}[0.022S_3 + 0.744C_3]\dot{q}_2^2 - [0.00125C_4S_5]\dot{q}_5^2 - [0.00125C_4S_5]\dot{q}_4^2 \\
&\quad + [-0.0025C_{23}C_4S_5 + 0.00164S_{23} + 0.00003S_{23}(1 - 2SS_4)]\dot{q}_1\dot{q}_4 \\
&\quad - [0.0025C_{23}S_4C_5 + 0.000642C_{23}S_4]\dot{q}_1\dot{q}_5 - [0.0025S_5]\dot{q}_2\dot{q}_5 \\
&\quad - [0.0025S_5]\dot{q}_3\dot{q}_5 - [0.0025S_4C_5]\dot{q}_4\dot{q}_5 - 8.4S_{23} + 0.25C_{23}
\end{align*}

\begin{align*}
n_4 &= \frac{1}{2}[0.0025SC_{23}S_4S_5 - 0.00086C_4S_5 + 0.00248C_2C_{23}S_4S_5]\dot{q}_1^2 \\
&\quad - \frac{1}{2}[0.00248C_3S_4S_5]\dot{q}_2^2 \\
&\quad + [0.00164S_{23} - 0.0025C_{23}C_4S_5 + 0.00248S_2C_4S_5 + 0.00003S_{23}(1 - 2SS_4)]\dot{q}_1\dot{q}_2 \\
&\quad + [0.0025C_{23}C_4C_5 - 0.00164S_{23} - 0.003S_{23}(1 - 2SS_4)]\dot{q}_1\dot{q}_3 \\
&\quad - [0.000642C_{23}S_4]\dot{q}_1\dot{q}_5 + [0.000642S_4]\dot{q}_2\dot{q}_5 \\
&\quad + [0.000642S_4]\dot{q}_3\dot{q}_5 + 0.028S_{23}S_4S_5
\end{align*}

\begin{align*}
n_5 &= \frac{1}{2}[0.0025(SS_{23}S_5 - SC_{23}C_4C_5) + 0.00248C_2(S_{23}S_5 - C_{23}C_4C_5) \\
&\quad - 0.000864S_4C_5]\dot{q}_1^2 - \frac{1}{2}[0.0025S_5 - 0.00248(C_3C_4C_5 + S_3S_5)]\dot{q}_2^2 \\
&\quad + [0.0026C_{23}S_4C_5 - 0.00248S_2S_4C_5 + 0.000642C_{23}S_4]\dot{q}_1\dot{q}_2 \\
&\quad + [0.0025C_{23}S_4C_5 + 0.000642C_{23}S_4]\dot{q}_1\dot{q}_3 + [0.000642C_{23}S_4]\dot{q}_1\dot{q}_4 \\
&\quad - [0.000642S_4]\dot{q}_2\dot{q}_4 - [0.000642S_4]\dot{q}_3\dot{q}_4 - 0.028(C_{23}C_5 + S_{23}C_4C_5)
\end{align*}

\begin{equation*}
n_6 = 0
\end{equation*}

\begin{figure}[H]
\centering
\fbox{\parbox{0.8\textwidth}{\centering 图 C.3.1：PUMA 560 机械臂}}
\caption{PUMA 560 机械臂示意图}
\end{figure}

\vspace{1em}

\subsection*{参考文献}

\noindent [Armstrong et al. 1986] Armstrong, B., O. Khatib, and J. Burdick, ``The explicit dynamic model and inertial parameters of the PUMA 560 arm,'' \textit{Proc. 1986 IEEE Conf. Robot. Autom.}, pp. 510--518, San Francisco, Apr. 7--10, 1986.

\vspace{0.5em}

\noindent [Bejczy 1974] Bejczy, A.K., ``Robot arm dynamics and control,'' NASA-JPL Technical Memorandum 33--669, 1974.

\vspace{0.5em}

\noindent [Paul 1981] Paul, R.P., \textit{Robot Manipulators: Mathematics, Programming and Control}. Cambridge, MA: MIT Press, 1981.

\ifstandalone
  \end{document}
\fi
